\section{Problem Definition and Physical Foundations}
\label{sec:foundations}

\subsection{Hamiltonian}
\label{subsec:hamiltonian}

The single-particle tight-binding Hamiltonian is constructed by projecting the continuous Schrödinger equation onto the discrete Wannier basis. 

In \textbf{Second Quantization} (the most common description in papers we read) with creation ($a_i^\dagger$) and annihilation ($a_i$) operators. It handles wave function anti-symmetrization, particle indistinguishability, and the Pauli exclusion principle. Because the tight-binding model preserves Time-Reversal Symmetry and SU(2) spin rotation symmetry without magnetic fields, the spin-up and spin-down transport channels are identical. We therefore neglect spin:
\begin{equation}
\mathcal{H} = \sum_{i} \epsilon_i a_i^\dagger a_i + \sum_{\langle i,j \rangle} t_{ij} \left( a_i^\dagger a_j + a_j^\dagger a_i \right)
\end{equation}
Here, $a_i^\dagger a_i$ is the number operator tracking site occupation, multiplied by the on-site potential $\epsilon_i$. The operator $a_i^\dagger a_j$ annihilates an electron at $j$ and creates one at $i$, driving the kinetic hopping with amplitude $t_{ij} = -t$.

In \textbf{First Quantization} (we use this for the simulations), the Hamiltonian is written using projection operators $|i\rangle\langle j|$, which physically represent an electron transitioning from the localized orbital at site $j$ to site $i$:
\begin{equation}
\mathcal{H} = \sum_{i} \epsilon_i |i\rangle\langle i| + \sum_{\langle i,j \rangle} t_{ij} \left( |i\rangle\langle j| + |j\rangle\langle i| \right)
\end{equation}

\textbf{Wannier Orbitals:} A complete, orthonormal set of basis functions $|i\rangle$ created via a unitary Fourier transform of the Bloch states. They are localized around the atomic lattice sites $\mathbf{R}_i$. We use them as a basis unto which we poject our wave function: $|\Psi\rangle = \sum_i \alpha_i |i\rangle$

\textbf{Bloch States:} fully delocalized eigenstates $\psi_{\mathbf{k}}(\mathbf{r})$ of a perfect periodic crystal. In our disordered lattice, the continuous translational Bloch symmetry is broken by randomness.

Because quantum percolation evaluates non-interacting electrons, restricting the problem to the single-particle Hilbert space allows the transfer matrices to scale linearly with the cross-sectional area ($\mathcal{O}(M^2)$). 
%Attempting to simulate the full second-quantized Hamiltonian would force operations within the many-particle Fock space, causing dimensions to scale exponentially, making large-scale computational evaluations impossible.

The disorder is introduced through the diagonal on-site potential energies $\epsilon_i$, sampled from an binary probability distribution representing the fraction $p$ of occupied sites:
\begin{equation}
\mathcal{P}(\epsilon_i) = p \delta(\epsilon_i) + (1-p) \delta(\epsilon_i - V_{\text{absent}})
\end{equation}

\subsection{Percolation}

To model site dilution (a lattice vacancy) without dynamically resizing the Hamiltonian matrices during computation, we assign missing sites an infinite potential barrier: $V_{\text{absent}} \to \infty$. This enforces the Dirichlet boundary condition $\psi_i = 0$ at the missing site's coordinate. 

In classical percolation theory, the threshold $p_c$ is determined by geometric connectivity: the existence of a continuous path through the lattice, or equivalently, an infinite cluster. 

The quantum percolation threshold $p_q$ is defined as the active site concentration where the absolute localization length $\xi$ diverges ($\xi \to \infty$). The threshold $p_q$ is determined via finite-size scaling of the dimensionless localization length $\Lambda_M=\xi_M/M$: the crossing point of $\Lambda_M$ vs. $p$ curves for different transverse system widths $M$ gives $p_q$. At $p<p_q$, all eigenstates are exponentially localized (finite $\xi$), whereas for $p>p_q$ extended states exist. One essentially looks at whether $\xi$ diverges or the dimensionless $\Lambda_M$ grows with $M$ (metallic behavior) at a given $p$, evaluating whether a quantum wavepacket can propagate macroscopically without its amplitude decaying to zero due to quantum interference.